\documentclass[11pt,a4paper]{article}
\usepackage[utf8]{inputenc}
\usepackage[french]{babel}
\usepackage[T1]{fontenc}
\usepackage{geometry}
\geometry{top=2.5cm, bottom=2.5cm, left=2.2cm, right=2.2cm}
\usepackage{xcolor}
\usepackage{hyperref}
\usepackage{listings}
\usepackage{tcolorbox}
\usepackage{tikz}
\usetikzlibrary{shapes.geometric, arrows.meta, positioning, shadows, fit}
\usepackage{array}
\usepackage{booktabs}
\usepackage{fancyhdr}
\usepackage{float}

% Palette de couleurs professionnelle
\definecolor{navyprimary}{RGB}{15, 23, 42}
\definecolor{blueaccent}{RGB}{29, 78, 216}
\definecolor{slategray}{RGB}{71, 85, 105}
\definecolor{codebg}{RGB}{248, 250, 252}
\definecolor{codeframe}{RGB}{226, 232, 240}
\definecolor{greenaccent}{RGB}{22, 101, 52}
\definecolor{purplaccent}{RGB}{126, 34, 206}
\definecolor{amberaccent}{RGB}{180, 83, 9}

% Configuration de hyperref
\hypersetup{
    colorlinks=true,
    linkcolor=blueaccent,
    filecolor=purplaccent,
    urlcolor=blueaccent,
    pdftitle={Rapport de Persistance des Données - IncidentFlow},
    pdfauthor={NetMar Technologies},
    pdfpagemode=UseOutlines
}

% Configuration de headheight pour fancyhdr
\setlength{\headheight}{14.5pt}
\pagestyle{fancy}
\fancyhf{}
\rhead{\textcolor{slategray}{\small IncidentFlow --- Persistance Utilisateurs, Rôles \& Workflows}}
\lhead{\textcolor{slategray}{\small Spécification Technique Backend}}
\rfoot{\textcolor{slategray}{\small Page \thepage}}
\lfoot{\textcolor{slategray}{\small NetMar Technologies}}
\renewcommand{\headrulewidth}{0.4pt}

% Définition de Java pour listings
\lstdefinelanguage{CustomJava}{
  language=Java,
  keywordstyle=\color{blueaccent}\bfseries,
  commentstyle=\color{slategray}\itshape,
  stringstyle=\color{greenaccent},
  numberstyle=\tiny\color{slategray},
  breaklines=true,
  showstringspaces=false,
  basicstyle=\ttfamily\small,
  tabsize=2
}

\lstset{
    backgroundcolor=\color{codebg},
    frame=single,
    rulecolor=\color{codeframe},
    captionpos=b,
    numbers=left,
    numbersep=8pt,
    xleftmargin=12pt
}

\begin{document}

% Page de titre
\begin{titlepage}
    \centering
    \vspace*{1.5cm}
    
    {\Huge \bfseries \color{navyprimary} IncidentFlow}\\[0.4cm]
    {\LARGE \color{blueaccent} Architectures \& Mécanismes de Persistance des Données}\\[0.2cm]
    {\large \color{slategray} Spécification Technique de Persistance JPA / Hibernate}\\[1.5cm]
    
    \rule{\linewidth}{0.5mm}\\[0.6cm]
    {\Large \bfseries Modélisation et Spécifications des Entités}\\[0.3cm]
    {\large \textbf{Utilisateurs, Rôles, Permissions \& Workflows Dynamiques}}\\[0.6cm]
    \rule{\linewidth}{0.5mm}\\[2cm]
    
    \begin{minipage}{0.45\textwidth}
        \begin{flushleft} \large
            \textbf{Auteur :}\\
            Équipe de Développement IncidentFlow\\
            NetMar Technologies
        \end{flushleft}
    \end{minipage}
    \begin{minipage}{0.45\textwidth}
        \begin{flushright} \large
            \textbf{Composants analysés :}\\
            Spring Data JPA / Hibernate 6\\
            PostgreSQL 16 \& Cache Redis
        \end{flushright}
    \end{minipage}

    \vfill
    {\large Année Universitaire / Stage 2026}
\end{titlepage}

\tableofcontents
\newpage

% 1. INTRODUCTION
\section{Introduction et Vue d'Ensemble de l'Architecture}

L'application \textbf{IncidentFlow} est une plateforme de gestion dynamique d'incidents IT et médicaux s'appuyant sur un moteur de \textit{State Machine} personnalisable. La robustesse et l'intégrité de la plateforme reposent sur la couche de persistance de données gérée par \textbf{Spring Data JPA} et \textbf{Hibernate ORM 6}, couplée à une base de données relationnelle \textbf{PostgreSQL 16} et à une couche de cache distribué \textbf{Redis}.

Ce document spécifie en détail la modélisation, le comportement transactionnel et les stratégies d'association JPA pour les trois piliers fondamentaux du système :
\begin{itemize}
    \item \textbf{La Gestion des Utilisateurs (\texttt{User})} : Authentification, profils et affectations organisationnelles.
    \item \textbf{La Matrice de Sécurité RBAC (\texttt{Role} \& \texttt{Permission})} : Attribution dynamique des droits système et contrôle d'accès fin.
    \item \textbf{Les Workflows Dynamiques (\texttt{Workflow}, \texttt{WorkflowState}, \texttt{WorkflowTransition})} : Définition visuelle des cycles de vie d'incidents, gestion des contraintes de transition et des rôles requis.
\end{itemize}

---

% 2. MODÈLE CONCEPTUEL ET SCHÉMA RELATIONNEL
\section{Modèle de Données et Relations JPA}

\subsection{Diagramme Entité-Association (ERD)}

Le schéma relationnel ci-dessous illustre la structure des tables PostgreSQL et les liens de clés étrangères reliant les utilisateurs, les rôles, les permissions et les composants de workflows.

\begin{figure}[H]
\centering
\begin{tikzpicture}[
    node distance=2.2cm and 2.8cm,
    entity/.style={rectangle, draw=navyprimary, fill=codebg, rounded corners, thick, minimum width=3.2cm, minimum height=1.2cm, align=center, font=\small\bfseries, drop shadow},
    arrow/.style={-{Stealth[scale=1.2]}, thick, draw=blueaccent}
]

% Entités Sécurité & Utilisateurs
\node[entity] (role) {Role\\ \normalfont\tiny (id, name, description)};
\node[entity, right=of role] (user) {User\\ \normalfont\tiny (id, email, password, role\_id)};
\node[entity, left=of role] (perm) {Permission\\ \normalfont\tiny (id, code, label, module)};

% Entités Workflows
\node[entity, below=2.5cm of role] (wf) {Workflow\\ \normalfont\tiny (id, name, category, version, active)};
\node[entity, below left=2cm and 0.8cm of wf] (state) {WorkflowState\\ \normalfont\tiny (id, name, label, workflow\_id)};
\node[entity, below right=2cm and 0.8cm of wf] (trans) {WorkflowTransition\\ \normalfont\tiny (id, from\_state, to\_state, role\_required)};

% Liens & Cardinalités
\draw[arrow] (perm) -- node[above, font=\tiny] {N:M (role\_permissions)} (role);
\draw[arrow] (role) -- node[above, font=\tiny] {1:N (role\_id)} (user);
\draw[arrow] (wf) -- node[left, font=\tiny] {1:N (Cascade.ALL)} (state);
\draw[arrow] (wf) -- node[right, font=\tiny] {1:N (Cascade.ALL)} (trans);

\end{tikzpicture}
\caption{Diagramme Entité-Association de la persistance IncidentFlow.}
\end{figure}

---

% 3. PERSISTANCE DES UTILISATEURS ET RÔLES (RBAC)
\section{Persistance des Utilisateurs, Rôles et Permissions (RBAC)}

\subsection{Entités \texttt{Role} et \texttt{Permission}}

Le système RBAC (\textit{Role-Based Access Control}) s'appuie sur une relation Many-to-Many entre les entités \texttt{Role} et \texttt{Permission}.

\begin{lstlisting}[language=CustomJava, caption={Mapping JPA de l'entite Role.}]
@Entity
@Table(name = "roles")
@Getter @Setter @NoArgsConstructor @AllArgsConstructor @Builder
public class Role implements Serializable {

    @Id
    @GeneratedValue(strategy = GenerationType.IDENTITY)
    private Long id;

    @Column(nullable = false, unique = true)
    private String name;

    private String description;

    @ManyToMany(fetch = FetchType.EAGER)
    @JoinTable(
        name = "role_permissions",
        joinColumns = @JoinColumn(name = "role_id"),
        inverseJoinColumns = @JoinColumn(name = "permission_id")
    )
    @Builder.Default
    private Set<Permission> permissions = new HashSet<>();
}
\end{lstlisting}

\subsubsection{Points Clés de Persistance RBAC :}
\begin{enumerate}
    \item \textbf{Chargement Immédiat (\texttt{FetchType.EAGER})} : Les permissions associées à un rôle sont chargées immédiatement lors de l'authentification pour alimenter le contexte de sécurité Spring Security sans requêtes N+1.
    \item \textbf{Jointure \texttt{role\_permissions}} : La table d'association conserve la correspondance exacte entre les identifiants de rôles et de privilèges.
\end{enumerate}

\subsection{Entité \texttt{User}}

L'utilisateur possède une association Many-to-One vers la table des rôles :

\begin{lstlisting}[language=CustomJava, caption={Extrait de l'entite User.}]
@Entity
@Table(name = "users")
public class User implements Serializable {

    @Id
    @GeneratedValue(strategy = GenerationType.IDENTITY)
    private Long id;

    @Column(nullable = false, unique = true)
    private String email;

    @Column(nullable = false)
    private String password;

    @ManyToOne(optional = false, fetch = FetchType.EAGER)
    @JoinColumn(name = "role_id", nullable = false)
    private Role role;
}
\end{lstlisting}

---

% 4. PERSISTANCE DES WORKFLOWS DYNAMIQUES
\section{Persistance des Workflows, États et Transitions}

La gestion des workflows dynamiques est la fonctionnalité cœur de l'application. Un workflow est composé d'un ensemble d'états (\texttt{WorkflowState}) et de règles de transition autorisées (\texttt{WorkflowTransition}).

\subsection{Entité \texttt{Workflow}}

L'entité principale \texttt{Workflow} agrège ses états et transitions via une cascade complète et une gestion stricte des entités orphelines (\texttt{orphanRemoval = true}).

\begin{lstlisting}[language=CustomJava, caption={Entite Workflow et ses collections associees.}]
@Entity
@Table(name = "workflows")
@Getter @Setter @NoArgsConstructor @AllArgsConstructor @Builder
public class Workflow implements Serializable {

    @Id
    @GeneratedValue(strategy = GenerationType.IDENTITY)
    private Long id;

    @Column(nullable = false)
    private String name;

    @Column(nullable = false)
    private String category;

    @Column(nullable = false)
    @Builder.Default
    private int version = 1;

    @Column(nullable = false)
    @Builder.Default
    private boolean active = true;

    @OneToMany(mappedBy = "workflow", cascade = CascadeType.ALL, orphanRemoval = true, fetch = FetchType.EAGER)
    @Builder.Default
    private List<WorkflowState> states = new ArrayList<>();

    @OneToMany(mappedBy = "workflow", cascade = CascadeType.ALL, orphanRemoval = true, fetch = FetchType.EAGER)
    @Builder.Default
    private List<WorkflowTransition> transitions = new ArrayList<>();
}
\end{lstlisting}

\subsection{Entités \texttt{WorkflowState} et \texttt{WorkflowTransition}}

Chaque transition contient désormais un champ \texttt{roleRequired} permettant de restreindre la réalisation d'un changement d'état à un rôle spécifique.

\begin{lstlisting}[language=CustomJava, caption={Entite WorkflowTransition avec restriction de role.}]
@Entity
@Table(name = "workflow_transitions")
@Getter @Setter @NoArgsConstructor @AllArgsConstructor @Builder
public class WorkflowTransition implements Serializable {

    @Id
    @GeneratedValue(strategy = GenerationType.IDENTITY)
    private Long id;

    @Column(name = "from_state", nullable = false)
    private String fromState;

    @Column(name = "to_state", nullable = false)
    private String toState;

    @Column(name = "role_required")
    private String roleRequired; // Exemple: "Administrateur", "Responsable", null

    @Column(name = "requires_comment", nullable = false)
    @Builder.Default
    private boolean requiresComment = false;

    @ManyToOne(optional = false)
    @JoinColumn(name = "workflow_id", nullable = false)
    @JsonIgnore
    private Workflow workflow;
}
\end{lstlisting}

---

% 5. GESTION DES CYCLES DE VIE JPA ET RÉSOLUTIONS D'ERREURS
\section{Gestion du Cycle de Vie JPA \& Résolution de l'Erreur \textit{Detached Entity}}

\subsection{Analyse du Problème de Persistance (\textit{Detached Entity Passed to Persist})}

Lors de la mise à jour d'un workflow existant depuis l'interface Web (Frontend React), le contrôleur reçoit un objet JSON désérialisé contenant des états et des transitions ayant des identifiants (\texttt{id}) déjà générés.

Lorsque l'on tente d'affecter directement cette liste à une entité managée par Hibernate :
\begin{enumerate}
    \item L'instance désérialisée n'appartient pas au contexte de persistance (\textit{Persistence Context}) courant.
    \item Hibernate détecte un identifiant non nul sur des objets enfants et interprète l'objet comme une entité détachée (\textit{Detached Entity}).
    \item L'appel à la méthode \texttt{persist()} au travers de la cascade \texttt{CascadeType.ALL} déclenche l'exception :\\
    \small\texttt{PersistentObjectException: detached entity passed to persist}.
\end{enumerate}

\subsection{Solution Implémentée dans \texttt{WorkflowService}}

Pour résoudre ce problème de manière définitive et garantir l'intégrité référentielle, le service procède par mise à jour en place sur la collection managée (\textit{In-Place Update}) ou par versioning automatique si des incidents sont déjà rattachés au workflow.

\begin{lstlisting}[language=CustomJava, caption={Algorithme de sauvegarde securise dans WorkflowService.java}]
@Transactional
@CacheEvict(value = {"workflows", "workflows-category"}, allEntries = true)
public Workflow saveWorkflow(Workflow workflow) {
    validateWorkflowGraph(workflow);

    if (workflow.getId() != null) {
        boolean hasLinkedIncidents = incidentRepository.existsByWorkflowId(workflow.getId());
        
        // Cas 1 : Des incidents sont lies -> Creation d'une nouvelle version (Versioning)
        if (hasLinkedIncidents) {
            Workflow original = workflowRepository.findById(workflow.getId())
                    .orElseThrow(() -> new ResourceNotFoundException("Workflow original non trouve"));
            original.setActive(false);
            workflowRepository.save(original);

            Workflow newVersion = new Workflow();
            newVersion.setName(workflow.getName());
            newVersion.setCategory(workflow.getCategory());
            newVersion.setVersion(original.getVersion() + 1);
            newVersion.setActive(true);

            // Re-instanciation propre sans ID pour eviter l'etat detache
            if (workflow.getStates() != null) {
                for (WorkflowState state : workflow.getStates()) {
                    newVersion.getStates().add(WorkflowState.builder()
                            .name(state.getName()).label(state.getLabel())
                            .colorClass(state.getColorClass()).active(state.isActive())
                            .workflow(newVersion).build());
                }
            }
            if (workflow.getTransitions() != null) {
                for (WorkflowTransition transition : workflow.getTransitions()) {
                    newVersion.getTransitions().add(WorkflowTransition.builder()
                            .fromState(transition.getFromState()).toState(transition.getToState())
                            .roleRequired(transition.getRoleRequired())
                            .requiresComment(transition.isRequiresComment())
                            .workflow(newVersion).build());
                }
            }
            return workflowRepository.save(newVersion);
        }

        // Cas 2 : Aucun incident lie -> Nettoyage et mise a jour en place
        Workflow existing = workflowRepository.findById(workflow.getId())
                .orElseGet(() -> workflow);

        if (existing != workflow) {
            existing.setName(workflow.getName());
            existing.setCategory(workflow.getCategory());
            existing.setActive(workflow.isActive());

            // Nettoyage de la collection managee (declenche orphanRemoval)
            existing.getStates().clear();
            if (workflow.getStates() != null) {
                for (WorkflowState state : workflow.getStates()) {
                    existing.getStates().add(WorkflowState.builder()
                            .name(state.getName()).label(state.getLabel())
                            .colorClass(state.getColorClass()).active(state.isActive())
                            .workflow(existing).build());
                }
            }

            existing.getTransitions().clear();
            if (workflow.getTransitions() != null) {
                for (WorkflowTransition transition : workflow.getTransitions()) {
                    existing.getTransitions().add(WorkflowTransition.builder()
                            .fromState(transition.getFromState()).toState(transition.getToState())
                            .roleRequired(transition.getRoleRequired())
                            .requiresComment(transition.isRequiresComment())
                            .workflow(existing).build());
                }
            }
            return workflowRepository.save(existing);
        }
    }
    return workflowRepository.save(workflow);
}
\end{lstlisting}

---

% 6. SYNCHRONISATION FRONTEND ET CYCLE DE VIE REST
\section{Flux de Persistance Frontend-Backend}

Le diagramme de séquence suivant résume le parcours de persistance lors de la modification du rôle autorisé d'une transition depuis l'interface utilisateur React.

\begin{figure}[H]
\centering
\begin{tikzpicture}[
    node distance=1.8cm,
    actor/.style={rectangle, draw=navyprimary, fill=codebg, rounded corners, font=\small\bfseries, minimum width=2.5cm, minimum height=0.8cm},
    line/.style={draw=slategray, dashed, thick}
]

\node[actor] (ui) {Interface React};
\node[actor, right=3.5cm of ui] (ctrl) {WorkflowController};
\node[actor, right=3.5cm of ctrl] (srv) {WorkflowService};
\node[actor, right=3.5cm of srv] (db) {PostgreSQL};

% Lignes de vie
\draw[line] (ui) -- ++(0,-5);
\draw[line] (ctrl) -- ++(0,-5);
\draw[line] (srv) -- ++(0,-5);
\draw[line] (db) -- ++(0,-5);

% Messages
\draw[-{Stealth}, thick, color=blueaccent] (ui.south) ++(0,-0.8) -- node[above, font=\tiny] {1. PUT /api/workflows/1 (JSON payload avec roleRequired)} ++(3.5,0);
\draw[-{Stealth}, thick, color=blueaccent] (ctrl.south) ++(0,-1.6) -- node[above, font=\tiny] {2. saveWorkflow(workflowDetails)} ++(3.5,0);
\draw[-{Stealth}, thick, color=blueaccent] (srv.south) ++(0,-2.4) -- node[above, font=\tiny] {3. In-Place update \& orphanRemoval} ++(3.5,0);
\draw[-{Stealth}, thick, color=greenaccent] (db.south) ++(0,-3.2) -- node[above, font=\tiny] {4. UPDATE workflows / DELETE \& INSERT transitions} ++(-3.5,0);
\draw[-{Stealth}, thick, color=greenaccent] (srv.south) ++(0,-4.0) -- node[above, font=\tiny] {5. Retour Workflow persistant (200 OK)} ++(-7,0);

\end{tikzpicture}
\caption{Diagramme de Séquence de la persistance d'une transition modifiée.}
\end{figure}

---

% 7. CONCLUSION ET BONNES PRATIQUES
\section{Conclusion et Bonnes Pratiques Retenues}

La stratégie de persistance mise en œuvre pour les entités \texttt{User}, \texttt{Role} et \texttt{Workflow} au sein d'IncidentFlow garantit :
\begin{itemize}
    \item \textbf{Intégrité des données} : Aucune incohérence d'état n'est possible grâce aux validations graphiques par parcours en profondeur (DFS) et au nettoyage automatique des sous-collections via \texttt{orphanRemoval}.
    \item \textbf{Gestion fluide des versions} : En présence d'incidents associés, l'historique n'est jamais corrompu car les anciens workflows sont archivés et une nouvelle version active est créée.
    \item \textbf{Sécurité granulaire RBAC} : La restriction dynamique des transitions par rôle (\texttt{roleRequired}) est directement enregistrée et validée au niveau de la couche service lors de chaque passage d'état d'un incident.
\end{itemize}

\end{document}
